% Source: https://artofproblemsolving.com/community/c3h1646384_ntalgebra_problem
% Stable identity: aops-topic-1646384-q01
\textbf{Form of $n$.} If the roots are integers $r>0$ and $s<0$, Vieta's formulas give $r+s=-1$ and $rs=-n$. So $s=-(r+1)$ and $n=r(r+1)$. Conversely, for every positive integer $r$, $x^2+x-r(r+1)=(x-r)(x+r+1)$ has the roots $r$ and $-(r+1)$. So we need the least $r\ge1$ with $720=2^4\cdot3^2\cdot5$ dividing $r(r+1)$.

\textbf{Divisibility.} The numbers $r$ and $r+1$ are coprime, so each of $16$, $9$, $5$ divides $r$ or divides $r+1$; that is, $r\equiv0$ or $-1$ modulo each of $16$, $9$ and $5$. By the Chinese remainder theorem the eight sign patterns give eight residues modulo $720$; the positive representatives below $720$ are
\[
80,\ 144,\ 224,\ 495,\ 575,\ 639,\ 719 \quad(\text{and } 720\equiv0).
\]
For example $r=80$: $80=16\cdot5$ and $81=9^2$. The least is $r=80$, giving
\[
n=80\cdot81=\boxed{6480}=9\cdot720 .
\]
Check: $x^2+x-6480=(x-80)(x+81)$.
